\section{Introdução}
\label{cap:introducao}

Os grandes modelos de linguagem (\textit{Large Language Models} -- LLMs)
permitem interações em linguagem natural, mas essa flexibilidade também pode
ser explorada para contornar suas restrições de segurança. Os ataques de
\textit{jailbreak} procuram induzir respostas que deveriam ser recusadas,
expondo limitações das salvaguardas desses sistemas
\cite{zou2023universaltransferableadversarialattacks,zeng2024johnny,li2025prefillleveljailbreakblackboxrisk}.
A diversidade de estratégias investigadas torna relevante compreender como
elas se relacionam e quais implicações apresentam para a segurança dos modelos.

Este artigo tem como objetivo analisar criticamente as abordagens de
jailbreak discutidas na literatura selecionada, identificando suas diferenças,
aproximações e possibilidades de complementaridade. A questão que orienta
a discussão é: como essas abordagens exploram diferentes dimensões da interação
com os LLMs e o que revelam sobre os limites de suas defesas?

Adota-se uma abordagem exploratória fundamentada em levantamento bibliográfico,
articulando os mecanismos e os resultados relatados nos estudos, sem validação
experimental própria. A análise privilegia as relações conceituais entre as
técnicas, considerando que diferenças nos protocolos de avaliação limitam
comparações diretas de eficácia. Com isso, busca-se oferecer uma compreensão
integrada das vulnerabilidades examinadas e de suas implicações defensivas.
